\documentclass[12pt,a4paper]{amsart}

\usepackage{amsmath,amssymb,amsthm}
\usepackage{mathtools}
\usepackage{enumitem}
\usepackage{hyperref}

\newtheorem{theorem}{Theorem}[section]
\newtheorem{lemma}[theorem]{Lemma}
\newtheorem{proposition}[theorem]{Proposition}
\newtheorem{corollary}[theorem]{Corollary}
\theoremstyle{definition}
\newtheorem{definition}[theorem]{Definition}
\newtheorem{assumption}[theorem]{Assumption}
\newtheorem{problem}[theorem]{Open Problem}
\newtheorem{conjecture}[theorem]{Conjecture}
\theoremstyle{remark}
\newtheorem{remark}[theorem]{Remark}

\newcommand{\Rig}{\textup{[Rig]}}
\newcommand{\Cond}{\textup{[Cond]}}
\newcommand{\Heur}{\textup{[Heur]}}
\newcommand{\Open}{\textup{[Open]}}

\begin{document}

\title[Lagrangian Singularity Transport under Fourier Duality]{%
  Lagrangian Singularity Transport under Fourier Duality:\\
  Cubic WKB Phases Induce A$_2$-Fold Singularities\\
  Whose Normal Form Is the Airy Distribution}

\author{[Author]}

\begin{abstract}
Paper~XV reduces the intermediate-regime
problem to whether the fold amplitude~$\Phi$
satisfies an estimate improving upon
Ces\`aro.
This paper establishes the microlocal
mechanism via three separate results.
\begin{enumerate}[\rm(I)]
\item \textit{(Local Airy normal form.)}
  Theorem~\ref{thm:airy-normalform}:
  Under the uniform A$_2$-stability
  conditions of Assumption~\ref{ass:microlocal},
  $\Phi(u)\sim e^{i\alpha_*u}a(u,c)
  \mathrm{Ai}(c_3^{-1/3}u)$
  in the Airy symbol class
  $a\in S^{-1/2}_{1,0}$,
  \emph{microlocally near the fold point};
  this is a local statement.
\item \textit{(Pointwise decay.)}
  Lemma~\ref{lem:pointwise}:
  $|\Phi(u)|\le C(1+|u|)^{-3/4}$
  uniformly for all $u$,
  combining the local normal form
  with an endpoint bound in the
  far zone.
\item \textit{(Global $L^1$-growth bound.)}
  Corollary~\ref{cor:global-bound}:
  $|\int_0^U\Phi\,du|=O(U^{1/4})$,
  from the pointwise bound.
  The Airy oscillation for $u<0$
  is accounted for entirely within
  the Airy normal form;
  no further cancellation is invoked.
\end{enumerate}
The forward direction is rigorous;
the converse requires
Hypothesis~\ref{hyp:FIO-rigidity}.
\end{abstract}

\maketitle

\section{Setup and Definitions}
\label{sec:setup}

\begin{definition}[A$_2$-fold singularity]
\label{def:A2-class}
Let $\mathcal{C}\subset T^*\mathbb{R}_\beta
\times T^*\mathbb{R}_u$ be a Lagrangian
submanifold.
Let $\pi_u:\mathcal{C}\to T^*\mathbb{R}_u$
denote the projection to the right factor.
\medskip

\noindent\textit{Geometric
(coordinate-free) definition.}
We say $\mathcal{C}$ has an
\emph{A$_2$-fold singularity at $p_0$} if:
\begin{enumerate}[\rm(i)]
\item $\det\,d\pi_u\big|_{p_0}=0$
  (rank-drop of $\pi_u$ by~1);
\item $d(\det\,d\pi_u)\big|_{p_0}\ne 0$
  (simple fold; no higher degeneracy).
\end{enumerate}
This is the A$_2$-type in Arnold's
classification \cite{Arnold1972}.
\medskip

\noindent\textit{Reduction to local
canonical coordinates.}
For a one-dimensional FIO
with phase $\Psi(\beta,u)
=\Theta(\beta)+u\beta$,
the canonical relation is
$\mathcal{C}=\{(\beta,\Theta'(\beta);\,
u,-\beta)\}$.
In these coordinates:
\begin{align}
  \label{eq:fold-i}
  \det\,d\pi_u=0
  \;&\iff\;
  \Theta''(\beta_0)=0,
  \\
  \label{eq:fold-ii}
  d(\det\,d\pi_u)\ne 0
  \;&\iff\;
  \Theta'''(\beta_0)\ne 0.
\end{align}
These are the coordinate expression
of~(i)--(ii) for this choice of phase;
the geometric conditions~(i)--(ii)
are invariant under canonical
transformations of $\mathcal{C}$.
\medskip

\noindent By the Thom--Boardman
classification \cite{Golub1973},
any one-dimensional FIO with an
A$_2$-fold singularity is,
\emph{microlocally near $p_0$},
equivalent to a scalar multiple of
$\mathrm{Ai}$ under a local
symplectic coordinate change.
This is a local-in-phase-space statement;
it does not extend globally.
\end{definition}

\begin{remark}[Polynomial phases
  are coordinate representations]
\label{rem:phases-coordinates}
The canonical invariant is the
Lagrangian submanifold
$\Lambda=\{(\beta,\Theta'(\beta))\}
\subset T^*\mathbb{R}$
and its Thom--Boardman stratum.
The expressions $e^{ic_3\beta^3}$
and $\mathrm{Ai}(c_3^{-1/3}u)$
are coordinate representations
of the same A$_2$ geometry
near the fold point;
this identification is local.
\end{remark}

\begin{remark}[Microlocal transport
  under Fourier duality]
\label{rem:microlocal-transport}
The Fourier transform acts on canonical
relations by the exchange
$(\beta,\xi;\,u,\eta)
\mapsto(u,\eta;\,\beta,\xi)$.
If the Lagrangian map germ of
$\pi_u:\mathcal{C}\to T^*\mathbb{R}_u$
belongs to the stable A$_2$-equivalence
class (Whitney fold, Mather-stable
\cite{Golub1973}),
then the germ of
$\pi_\beta:\mathcal{C}\to
T^*\mathbb{R}_\beta$
also belongs to the stable A$_2$-class.
This is a \emph{local statement about
stable equivalence of Lagrangian map
germs}; it fails if the fold is
not Mather-stable ($c_3(c)\to 0$).
\end{remark}

\subsection{The fold amplitude}

\begin{equation}\label{eq:Phi-def}
  \Phi(u)
  :=\int_{\mathbb{R}}
  |\beta|^{-1/2}\chi(\alpha_*+\beta)
  e^{i\beta u}\,d\beta,
  \quad u\in\mathbb{R},
\end{equation}
$\chi\in C^\infty_c$,
$\alpha_*>0$ fold point.
$g(\alpha)
:= |\alpha-\alpha_*|^{-1/2}
e^{i\theta(\alpha)}\chi(\alpha)$,
$\theta''(\alpha_*)=0$,
$c_3:=\frac{1}{6}\theta'''(\alpha_*)\ne 0$.
$\beta=\alpha-\alpha_*$,
$\Theta(\beta):=\theta(\alpha_*+\beta)$:
\begin{equation}\label{eq:Phi-WKB}
  \Phi(u)
  = e^{i\alpha_* u}
  \int_{\mathbb{R}}
  |\beta|^{-1/2}
  e^{i(\Theta(\beta)+\beta u)}
  \chi(\alpha_*+\beta)
  \,d\beta.
\end{equation}

\begin{theorem}[Endpoint decay \Rig]
\label{thm:endpoint}
$|\Phi(u)|=O(|u|^{-1/2})$
for all $\theta\in C^\infty$.
\end{theorem}

\begin{proof}
van~der~Corput with weight
$|\beta|^{-1/2}$; Erd\'elyi
for $\theta\equiv 0$.
\end{proof}

\section{Cubic Phase Induces A$_2$-Fold}
\label{sec:nonpreservation}

\begin{assumption}[Local stable A$_2$-fold \Cond]
\label{ass:fold-local}
$\Theta(\beta)=c_3\beta^3+O(\beta^4)$,
$c_3\ne 0$, and:
\begin{enumerate}[\rm(F1)]
\item $\Theta''(\beta)\ne 0$
  for $\beta\ne 0$ near $0$;
  at most one fold point, no cusp.
\item $\int|\beta|^{-1/2}
  e^{i\Theta}\chi\,d\beta
  \not\equiv 0$ near $0$.
\item $\Theta'$ has exactly one
  critical point ($\beta=0$)
  on $\mathrm{supp}\,\chi$.
\end{enumerate}
\end{assumption}

\begin{proposition}[Cubic phase
  $\Rightarrow$ local Airy normal form \Rig]
\label{prop:nonpreservation}
Under Assumption~\ref{ass:fold-local},
let $F(u):=\int|\beta|^{-1/2}
e^{i(\Theta(\beta)+\beta u)}
\chi(\beta)\,d\beta$.
\begin{enumerate}[\rm(1)]
\item \textit{(A$_2$-fold.)} $\mathcal{C}$
  has an A$_2$-fold singularity at $(0,0)$
  (by~\eqref{eq:fold-i}--\eqref{eq:fold-ii}
  and (F1)).
\item \textit{(Local CFU normal form.)} 
  Fix a fold neighbourhood
  $U_\epsilon=\{|\beta|<\epsilon\}$
  and a far region
  $V=\mathrm{supp}\,\chi\setminus
  U_{\epsilon/2}$.
  Let $1=\phi_{\mathrm{fold}}
  +\phi_{\mathrm{far}}$
  be a smooth partition subordinate
  to $\{U_\epsilon, V\}$.
  \begin{enumerate}[\rm(a)]
  \item \textit{Fold part.}
    In $U_\epsilon$,
    the CFU theorem \cite{CFU1957}
    provides a smooth local
    coordinate change
    $\beta\mapsto t(\beta,u)$,
    a local diffeomorphism
    on $U_\epsilon$ (not global),
    and smooth functions $R(u)$,
    $\mu(u)=c_3^{-1/3}u+O(u^2)$
    such that:
    \begin{equation}\label{eq:CFU-nf}
      \Theta(\beta)+u\beta
      \sim\tfrac{1}{3}t(\beta,u)^3
      - \mu(u)\,t(\beta,u)
      + R(u)
      +O(\beta^4)
    \end{equation}
    as $\beta\to 0$, uniformly in $u$.
    This is an \emph{asymptotic
    generating function equivalence}
    in $U_\epsilon$;
    it is not an exact global identity.
    The Jacobi factor
    $J(t,u):=d\beta/dt$
    satisfies
    $J(t,u)=c_3^{-1/6}|u|^{-1/2}
    (1+O(|u|^{-1}))$
    for large $|u|$
    (standard CFU expansion
    \cite{CFU1957,Olver1974}).
  \item \textit{Far part.}
    On $V$, $\Theta''\ne 0$ by (F1),
    so $F_{\mathrm{far}}(u)
    :=\int\phi_{\mathrm{far}}
    |\beta|^{-1/2}
    e^{i(\Theta+u\beta)}\,d\beta$
    is a non-degenerate oscillatory
    integral;
    van~der~Corput gives
    $|F_{\mathrm{far}}(u)|
    \le C|u|^{-1/2}$.
  \end{enumerate}
\item \textit{(Symbol class
  $S^{-1/2}_{1,0}$ in $u$:)
  symbol tracking.}
  In the fold region,
  the amplitude after CFU is:
  $|\beta(t,u)|^{-1/2}J(t,u)$.
  By H\"ormander's parametric
  stationary phase theorem
  \cite{Hormander1983}, Ch.~7.7:
  integrating over $t$ with smooth
  amplitude against the Airy kernel
  $e^{i(t^3/3-\mu t)}$
  yields a symbol of order $-1/2$
  in $u$,
  \emph{provided the following
  scaling conditions hold:}
  \begin{enumerate}[\rm(S1)]
  \item The phase $t^3/3-\mu(u)t$
    is non-degenerate in $t$
    for $|t|\gg|\mu|^{1/2}$;
  \item $|\mu(u)|
    \sim c_3^{-1/3}|u|\to\infty$
    as $|u|\to\infty$
    (from Assumption~\ref{ass:microlocal}(ii));
  \item $J(t,u)
    \sim c_3^{-1/6}|u|^{-1/2}$
    uniformly
    (Assumption~\ref{ass:microlocal}(iv)).
  \end{enumerate}
  Under (S1)--(S3),
  $a_0(u,c):=e^{-iR(u)}
  \int_{\mathbb{R}}
  |\beta(t,u)|^{-1/2}J(t,u)\,
  e^{i(t^3/3-\mu(u)t)}\,dt
  \big/\mathrm{Ai}(-\mu(u))$
  satisfies
  $a_0\in S^{-1/2}_{1,0}$
  in the $u$-variable.
  This is \emph{not} directly
  inherited from $|\beta|^{-1/2}$
  on the $\beta$-side;
  it arises from the Jacobian
  $J\sim c_3^{-1/6}|u|^{-1/2}$
  under the CFU scaling.
\end{enumerate}
\end{proposition}

\section{Normal Form, Decay,
  and Global Integrability}
\label{sec:normalform}

\begin{assumption}[Uniform A$_2$-stability \Cond]
\label{ass:microlocal}
Assumption~\ref{ass:fold-local}
holds uniformly in $c\ge c_0$.
In addition, uniformly in $c$:
\begin{enumerate}[\rm(i)]
\item $\Theta(\beta)=c_3\beta^3+O(\beta^4)$
  in a $c$-independent neighbourhood;
\item $c_3(c)\ge\delta_0>0$;
\item $|\Theta''(\beta^*(u))|
  \ge\delta_1|u|^{1/2}$
  for $\delta_2\le|u|\le c^{\theta_0/3}$;
\item $0<\delta_3\le|J(t,u)|
  \le\delta_3^{-1}$;
\item (F3) holds uniformly:
  no secondary stationary points
  for all $c\ge c_0$.
\end{enumerate}
\end{assumption}

\begin{theorem}[Local Airy normal form \Rig]
\label{thm:airy-normalform}
Under Assumption~\ref{ass:microlocal},
\emph{microlocally near the fold point
$(\beta_0,u_0)=(0,0)$
and for $\delta_2\le|u|\le c^{\theta_0/3}$}:
\begin{equation}\label{eq:airy-normalform}
  \Phi(u)
  \;\sim\;
  e^{i\alpha_* u}\,
  a(u,c)\,
  \mathrm{Ai}\!\bigl(c_3^{-1/3}u\bigr),
\end{equation}
where $a(u,c)\in S^{-1/2}_{1,0}$
in the $u$-variable under
scaling conditions (S1)--(S3)
of Proposition~\ref{prop:nonpreservation}(3):
$|\partial_u^j a|
\le C_j(1+|u|)^{-1/2-j}$
uniformly in $c$.
The symbol $\sim$ in~\eqref{eq:airy-normalform}
denotes asymptotic equivalence
in the fold zone;
the far-zone contribution
$F_{\mathrm{far}}$ is bounded
separately by
Theorem~\ref{thm:endpoint}.
\end{theorem}

\begin{proof}[Proof sketch]
Apply the local CFU
coordinate change of
Proposition~\ref{prop:nonpreservation}(2a)
to the fold part in $U_\epsilon$;
conditions (iii)--(v) ensure
non-degeneracy, Jacobian bounds,
and no secondary stationary points.
The far part is bounded by
van~der~Corput.
The symbol class follows from
H\"ormander's parametric stationary
phase under (S1)--(S3).
\end{proof}

\begin{lemma}[Pointwise decay \Rig]
\label{lem:pointwise}
Under Assumption~\ref{ass:microlocal}:
\begin{equation}\label{eq:pointwise}
  |\Phi(u)|\le C(1+|u|)^{-3/4}
  \quad\text{uniformly for all }u.
\end{equation}
\end{lemma}

\begin{proof}
Decompose via
$\Phi=\Phi_{\mathrm{fold}}
+\Phi_{\mathrm{far}}$.
\medskip

\noindent\textit{Fold zone
($|u|\le 2c^{\theta_0/3}$).}
By Theorem~\ref{thm:airy-normalform}:
$|\Phi_{\mathrm{fold}}(u)|
\lesssim|a(u,c)|\cdot
|\mathrm{Ai}(c_3^{-1/3}u)|$.
\begin{itemize}
\item $u\ge 0$:
  $|\mathrm{Ai}(c_3^{-1/3}u)|
  \le Cu^{-1/4}e^{-c'u^{3/2}}$
  \cite{Olver1974},
  $|a|\le Cu^{-1/2}$:
  product $\le Cu^{-3/4}$.
\item $u<0$:
  $|\mathrm{Ai}(c_3^{-1/3}u)|
  \le C|u|^{-1/4}$
  \cite{Olver1974},
  $|a|\le C|u|^{-1/2}$:
  product $\le C|u|^{-3/4}$.
\end{itemize}
Hence $|\Phi_{\mathrm{fold}}|
\le C(1+|u|)^{-3/4}$
in the fold zone.
\medskip

\noindent\textit{Far zone
($|u|\ge c^{\theta_0/3}$).}
Theorem~\ref{thm:endpoint}:
$|\Phi_{\mathrm{far}}(u)|
\le C|u|^{-1/2}
\le C|u|^{-3/4}$
(since $|u|\ge 1$).
\medskip

\noindent\textit{Transition.}
In the overlap
$c^{\theta_0/3}\le|u|
\le 2c^{\theta_0/3}$,
both bounds hold;
uniformity follows from
conditions (iii)--(iv).
\end{proof}

\begin{lemma}[$L^1$-growth estimate \Rig]
\label{lem:airy-integral}
Under Lemma~\ref{lem:pointwise}:
\begin{equation}\label{eq:L1-bound}
  \left|\int_0^U\Phi(u)\,du\right|
  \le C\int_0^U(1+u)^{-3/4}\,du
  =O(U^{1/4}).
\end{equation}
This bound uses only the pointwise
estimate~\eqref{eq:pointwise}.
The Airy oscillation for $u<0$
(which produces the $|u|^{-1/4}$ decay
via stationary phase within the
Airy representation)
is accounted for entirely within
the normal form of
Theorem~\ref{thm:airy-normalform};
no further oscillatory cancellation
is invoked at the level of the
$u$-integral.
\end{lemma}

\begin{corollary}[Ces\`aro improvement \Rig]
\label{cor:global-bound}
Under Assumption~\ref{ass:microlocal}:
$|\int_0^U\Phi\,du|=O(U^{1/4})$,
improving upon
Ces\`aro $O(U^{1/2})$.
The improvement is due to the
structural decay $(1+|u|)^{-3/4}$,
which combines:
(a)~$a\in S^{-1/2}_{1,0}$
under the CFU scaling
(Proposition~\ref{prop:nonpreservation}(3)),
(b)~$|\mathrm{Ai}(-t)|=O(t^{-1/4})$,
(c)~uniform transition
between fold and far zones.
\end{corollary}

\subsection{Converse and FIO rigidity}

\begin{hypothesis}[FIO rigidity \Cond]
\label{hyp:FIO-rigidity}
Suppose $\Phi$ is in the Airy class and:
\begin{enumerate}[\rm(a)]
\item $\sigma(g\mapsto\Phi)
  \big|_\mathcal{C}\ne 0$;
\item $g\mapsto\Phi$ is injective
  modulo smoothing;
\item $\mathrm{WF}(\Phi)
  =\pi_u(\mathcal{C})$ on the
  A$_2$ stratum.
\end{enumerate}
Then $\mathcal{C}$ contains an
A$_2$-fold uniformly in $c$.
\end{hypothesis}

\begin{remark}
Conditions (a)--(c) are standard
microlocal ellipticity conditions
\cite{Zworski2012}.
Verification for PSWF is Paper~XVII.
\end{remark}

\begin{corollary}[Conditional equivalence
  \Cond]
\label{cor:conditional-equiv}
Under Hypothesis~\ref{hyp:FIO-rigidity}:
Airy-class bound
$\iff$
Assumption~\ref{ass:microlocal}.
Without it: $\Rightarrow$ only.
\end{corollary}

\section{Residual Open Problem}
\label{sec:open}

\begin{problem}[A$_2$-stability
  for PSWF \Open]
\label{prob:PSWF-microlocal}
Does the PSWF amplitude satisfy
Assumptions~\ref{ass:fold-local}
and~\ref{ass:microlocal}
uniformly in $c\to\infty$?
\begin{enumerate}[\rm(a)]
\item $c_3(c)\ge\delta_0>0$?
  (Olver WKB \cite{Olver1974}.)
\item $\theta''(\alpha)\ne 0$
  for $\alpha\ne\alpha_*$, uniformly?
\item Uniform Jacobian
  Assumption~\ref{ass:microlocal}(iv)?
\end{enumerate}
Outcomes:
(A)~Assumptions hold:
$|\int_0^U\Phi|=O(U^{1/4})$.
(B)~$c_3(c)\to 0$:
Ces\`aro sharp.
\end{problem}

\begin{thebibliography}{9}
\bibitem{PaperXV}
[Author], \textit{Paper~XV:
Intermediate-Regime Kernel Estimates},
2026.

\bibitem{Erdelyi1956}
A.~Erd\'elyi,
\textit{Asymptotic Expansions},
Dover, 1956.

\bibitem{Olver1974}
F.~W.~J.~Olver,
\textit{Asymptotics and Special
Functions},
Academic Press, 1974.

\bibitem{Zworski2012}
M.~Zworski,
\textit{Semiclassical Analysis},
AMS, 2012.

\bibitem{Hormander1983}
L.~H\"ormander,
\textit{The Analysis of Linear PDE~I},
Springer, 1983.

\bibitem{CFU1957}
C.~Chester, B.~Friedman, F.~Ursell,
\textit{An extension of the method
of steepest descents},
Proc.\ Cambridge Philos.\ Soc.\
\textbf{53} (1957), 599--611.

\bibitem{Arnold1972}
V.~I.~Arnold,
\textit{Normal forms for functions
near degenerate critical points},
Funct.\ Anal.\ Appl.\
\textbf{6} (1972), 254--272.

\bibitem{Golub1973}
M.~Golubitsky, V.~Guillemin,
\textit{Stable Mappings and
Their Singularities},
Springer, 1973.
\end{thebibliography}

\end{document}
